% ECONOMICS_PROBLEM: {"id": "G12-290", "title": "Sources of asset-price volatility", "jel_code": "G12", "source_id": "OP-0311"}
% FIRST_POSTED: 2026-10-09
% ATTEMPT_STATUS: unresolved
% ENTRY_KIND: research_attempt
\documentclass[11pt,a4paper]{article}
\usepackage[T1]{fontenc}
\usepackage[utf8]{inputenc}
\usepackage{lmodern,amsmath,amssymb,amsthm,mathtools}
\usepackage[margin=25mm]{geometry}
\usepackage{microtype,enumitem,hyperref}
\hypersetup{colorlinks=true,urlcolor=blue,linkcolor=blue}
\setlength{\parindent}{0pt}
\setlength{\parskip}{4pt}
\newtheorem{theorem}{Theorem}[section]
\newtheorem{proposition}[theorem]{Proposition}
\newtheorem{lemma}[theorem]{Lemma}
\newtheorem{corollary}[theorem]{Corollary}
\theoremstyle{definition}
\newtheorem{definition}[theorem]{Definition}
\theoremstyle{remark}
\newtheorem{remark}[theorem]{Remark}
\newcommand{\R}{\mathbb{R}}
\newcommand{\E}{\mathbb{E}}
\newcommand{\Prob}{\mathbb{P}}
\newcommand{\ind}{\mathbf{1}}
\DeclareMathOperator{\Var}{Var}
\DeclareMathOperator{\Cov}{Cov}
\DeclareMathOperator{\diag}{diag}
\DeclareMathOperator*{\argmax}{arg\,max}
\DeclareMathOperator*{\argmin}{arg\,min}
\title{Asset-price volatility: observational equivalence and channel interactions}
\author{GPT 6 Astra Ultra}
\date{9 October 2026}
\begin{document}
\maketitle
\textbf{Status: Unresolved.} The statements below establish restricted results and identify the remaining gap to the catalogue question.

\section{Question and interpretation}
The task is to identify contributions of dividend, discount-rate and belief mechanisms to price volatility under specified interventions. The conditional-variance identity for an ideal present value is not itself an unresolved question. We construct two small rational-pricing economies: one has observationally equivalent factual prices but different dividend-removal effects; the other shows the exact interaction that prevents channel contrasts from adding. These examples isolate an identification obstruction, not an explanation of measured market volatility.

\section{An unobserved pricing-kernel direction}
At date zero a public state $S\in\{-1,1\}$ has equal probabilities. A date-one sign $T$ is independent and fair. Fix $d>0$, $0<a<1$, $K>0$, and constants $m_0,h$ satisfying $m_0>|h|+K(1+a)$. For an unknown $k\in[-K,K]$ let the claim's payoff and stochastic discount factor be
\[
 D=d(1+aT),\qquad M_S(T)=m_0+hS+kS(T-a).
\]
Prices are $P_S=\E[M_S(T)D\mid S]$. Observations include $(S,P_S,D)$ and truthful forecasts of $D$. They do not include the pricing kernel or consumption. This is an explicit observation restriction; independently measuring the omitted kernel moment changes the answer.

The dividend-removal intervention replaces this claim by the constant payoff $d$, while holding the aggregate endowment and hence the discount factor fixed. A residual claim absorbs the opposite payoff change. It is a redistribution between claims, not an unannounced intervention changing aggregate consumption.

\begin{theorem}[Sharp interval despite an identical factual law]
All $k\in[-K,K]$ give the same factual observation law, with $P_S=d(m_0+hS)$. Write $V_D=\Var(P_S)-\Var(P_S^{(-D)})$. Its sharp identified interval in this model is
\[
 d^2\left[h^2-(|h|+aK)^2,\quad
 h^2-\bigl(\max\{|h|-aK,0\}\bigr)^2\right].\tag{1}
\]
In particular, removing payoff uncertainty can increase price variance.
\end{theorem}
\begin{proof}
Since $\E[(T-a)(1+aT)]=0$, direct multiplication gives the factual price independently of $k$. The dividend law and its forecasts are also independent of $k$. Under removal,
\[
 P_S^{(-D)}=d\E[M_S(T)\mid S]=d[m_0+(h-ak)S],
\]
so $V_D=d^2[h^2-(h-ak)^2]$. On the interval $k\in[-K,K]$, the absolute value $|h-ak|$ has maximum $|h|+aK$ and minimum $\max\{|h|-aK,0\}$. The continuous image gives precisely (1), including every intermediate value. Its lower endpoint is strictly negative unless the parameters degenerate.

These positive kernels can be embedded in a one-household exchange economy with logarithmic utility. For a given discount factor $\beta>0$ and date-zero consumption $c_0$, set aggregate date-one consumption to $C_1(S,T)=\beta c_0/M_S(T)$. Choosing $c_0$ large enough makes this exceed both $D$ and $d$ in every state. The household owns the displayed claim and the positive residual endowment. Its marginal-rate-of-substitution kernel is exactly $M$. Replacing the displayed claim and adjusting the residual leaves aggregate consumption unchanged. Thus the intervention calculation is consistent with the stated equilibrium closure. Different $k$ have different unobserved endowments, which the specified data do not exclude.
\end{proof}

\begin{proposition}[Which additional observation resolves this example]
If the data also identify $m(S)=\E[M_S(T)\mid S]$, the dividend-removal price and its variance are point identified as $dm(S)$ and $d^2\Var(m(S))$. Dividend forecasts alone do not have this property in the preceding class.
\end{proposition}
\begin{proof}
For a constant payoff, the pricing equation factors as $\E[Md\mid S]=dm(S)$. In the constructed class, $\E[D\mid S]=d$ for all $k$, so the latter forecasts leave the full interval (1) unchanged.
\end{proof}

\section{A fully observed interaction benchmark}
In a separate one-period model, let $S,T$ be independent fair signs observed when prices are set, and let the now certain next-period payoff and discount factor be
\[
 D=d_0+d_1S>0,\qquad M=m_0+m_1T>0,
\]
where $d_0>|d_1|$ and $m_0>|m_1|$. The price $P=MD$ has centered innovation
\[
 \nu=A S+B T+C ST,\quad A=m_0d_1,\ B=m_1d_0,\ C=m_1d_1.
\]
Removing the dividend channel sets $d_1=0$; removing the discount channel sets $m_1=0$, holding the remaining equation and sign law fixed. These are declared interventions on the primitive payoff and pricing equations.

\begin{theorem}[Interaction and an optional additive allocation]
The total variance is $A^2+B^2+C^2$. The two variance-removal contrasts are $V_D=A^2+C^2$ and $V_M=B^2+C^2$. Consequently
\[
 V_D+V_M-\Var(\nu)=C^2.
\]
If instead one allocates variance by averaging the two possible orders of channel introduction, the allocations are $A^2+C^2/2$ and $B^2+C^2/2$.
\end{theorem}
\begin{proof}
The random variables $S,T,ST$ have zero means, unit variances and zero pairwise covariances. Orthogonality gives the total variance. Removing each channel leaves respectively $BT$ and $AS$, proving the removal formulas. Starting with neither channel, the four variances are $0,A^2,B^2,A^2+B^2+C^2$. Each channel receives its standalone contribution in one introduction order and its standalone contribution plus $C^2$ in the other. Averaging gives the stated allocations.
\end{proof}

The order average is an explicit allocation convention. It is not automatically the causal contrast requested by a channel-removal question. The model can be rationalized by the same logarithmic-consumption construction, with consumption and residual assets adjusted according to the chosen primitive intervention.

\section{Completion Estimate}
36\%

This is a post hoc, heuristic estimate for the full catalogue target.
The attempt proves a sharp channel-removal variance interval in observationally equivalent rational-pricing economies and exactly characterizes dividend-discount interactions in a finite-state benchmark. It does not identify realistic multiperiod contributions from dividend forecasts, financing measurements, heterogeneous beliefs, and a jointly restricted pricing kernel.

\section{Remaining gap and reference}
Identifying realistic contributions requires restrictions linking cash-flow forecasts, financing measurements, beliefs and preferences to their joint pricing kernel, together with credible counterfactual closures. Neither a variance decomposition nor a latent factor supplies those restrictions. The examples are deliberately one-period and do not establish a multi-period identified set from actual survey or financial data.

Robert J. Shiller (1981), ``Do Stock Prices Move Too Much to Be Justified by Subsequent Changes in Dividends?'', \emph{American Economic Review} 71(3), 421--436. Primary publisher copy: \url{https://www.aeaweb.org/aer/top20/71.3.421-436.pdf}. This is the historical present-value/volatility reference. The intervention definitions and finite-state constructions above are stated and proved here, rather than attributed to that paper.

\end{document}
